% Part: first-order-logic
% Chapter: models-theories
% Section: theories

\documentclass[../../../include/open-logic-section]{subfiles}

\begin{document}

\olfileid[hi]{fol}{mat}{the}

\olsection{प्रथम-क्रम सिद्धांतों के उदाहरण}

\begin{ex}
भाषा~$\Lang L_<$ में प्रबल रैखिक क्रमों का सिद्धांत निम्न समुच्चय द्वारा
अभिगृहीतीकृत है:
\begin{align*}
\{\quad & \lforall[x][\lnot x < x], \\
& \lforall[x][\lforall[y][((x < y \lor y <
    x) \lor x = y)]], \\
& \lforall[x][\lforall[y][\lforall[z][((x < y
      \land y < z) \lif x < z)]]] \quad \}
\end{align*}
अभीष्ट !!{structure}s इस अभिगृहीत-तंत्र द्वारा पूर्णतः निरूपित होती हैं: प्रत्येक प्रबल रैखिक
क्रम इस अभिगृहीत-तंत्र का प्रतिरूप है; और विलोमतः, यदि~$R$ किसी समुच्चय~$X$
पर रैखिक क्रम है, तो संरचना~$\Struct M$, जिसके लिए $\Domain M = X$ और
$\Assign{<}{M} = R$, इस सिद्धांत का प्रतिरूप है।
\end{ex}

\begin{ex}
$\Obj 1$ (!!{constant}) और~$\cdot$ (द्विस्थानी !!{function}) वाली भाषा में
समूहों का सिद्धांत निम्न वाक्यों द्वारा अभिगृहीतीकृत है:
\begin{align*}
& \lforall[x][\eq[(x \cdot \Obj 1)][x]]\\
& \lforall[x][\lforall[y][\lforall[z][\eq[(x \cdot (y \cdot z))][((x
          \cdot y) \cdot z)]]]]\\
& \lforall[x][\lexists[y][\eq[(x \cdot y)][\Obj 1]]]
\end{align*}
\end{ex}

\begin{ex}
पियानो अंकगणित का सिद्धांत अंकगणित की भाषा~$\Lang L_A$ के निम्न
वाक्यों द्वारा अभिगृहीतीकृत है:
\begin{align*}
& \lforall[x][\lforall[y][(\eq[x'][y'] \lif \eq[x][y])]]\\
& \lforall[x][\eq/[\Obj 0][x']]\\
& \lforall[x][\eq[(x + \Obj 0)][x]]\\
& \lforall[x][\lforall[y][\eq[(x + y')][(x + y)']]]\\
& \lforall[x][\eq[(x \times \Obj 0)][\Obj 0]]\\
& \lforall[x][\lforall[y][\eq[(x \times y')][((x \times y) + x)]]]\\
& \lforall[x][\lforall[y][(x < y \liff \lexists[z][\eq[(z' + x)][y])]]]\\
\intertext{और इस रूप के सभी वाक्य:}
& (!A(\Obj 0) \land \lforall[x][(!A(x) \lif !A(x'))]) \lif \lforall[x][!A(x)]
\end{align*}
चूँकि अंतिम रूप के अनंततः अनेक वाक्य हैं, यह अभिगृहीत-तंत्र अनंत है। अंतिम
रूप को \emph{आगमन-योजना} कहते हैं। (वास्तव में आगमन-योजना यहाँ बताए गए रूप
से थोड़ी अधिक जटिल है।)

अंतिम अभिगृहीत~$<$ की \emph{स्पष्ट परिभाषा} है।
\end{ex}

\begin{ex}
शुद्ध समुच्चयों का सिद्धांत गणित की नींव (और उसके दर्शन) में महत्त्वपूर्ण
भूमिका निभाता है। कोई समुच्चय शुद्ध है यदि उसके सभी !!{element}s भी शुद्ध
समुच्चय हों। इसलिए रिक्त समुच्चय शुद्ध माना जाता है, किंतु जिस समुच्चय का कोई
ऐसा !!a{element} हो जो स्वयं समुच्चय न हो, वह शुद्ध नहीं होगा। अतः शुद्ध
समुच्चय वे हैं जो केवल रिक्त समुच्चय से, और किसी ``मूल-अवयव''---अर्थात् ऐसी
वस्तु जो स्वयं समुच्चय नहीं---के बिना बनते हैं।

निम्न को शुद्ध समुच्चयों के सिद्धांत का अभिगृहीत-तंत्र माना जा सकता है:
\begin{align*}
& \lexists[x][\lnot \lexists[y][y \in x]]\\
& \lforall[x][\lforall[y][(\lforall[z](z \in x \liff z \in y) \lif 
\eq[x][y])]]\\
& \lforall[x][\lforall[y][\lexists[z][\lforall[u][(u \in z \liff
(\eq[u][x] \lor \eq[u][y]))]]]]\\
& \lforall[x][\lexists[y][\lforall[z][(z \in y \liff \lexists[u][(z \in
        u \land u \in x)])]]]\\
\intertext{और इस रूप के सभी वाक्य:} &
\lexists[x][\lforall[y][(y \in x \liff !A(y))]]
\end{align*}
पहला अभिगृहीत कहता है कि ऐसा समुच्चय है जिसका कोई !!{element} नहीं (अर्थात्
$\emptyset$ विद्यमान है); दूसरा कहता है कि समुच्चय विस्तारात्मक हैं; तीसरा
कहता है कि किन्हीं समुच्चयों~$X$ और~$Y$ के लिए समुच्चय~$\{X, Y\}$ विद्यमान
है; और चौथा कहता है कि किसी भी समुच्चय~$X$ के लिए समुच्चय~$\cup X$ विद्यमान
है, जहाँ~$\cup X$,~$X$ के सभी अवयवों का सम्मिलन है।

अंत में उल्लिखित !!{sentence}s को सामूहिक रूप से \emph{अनौपचारिक
समुच्चय-निर्माण योजना} कहते हैं। यह मूलतः कहती है कि प्रत्येक~$!A(x)$ के लिए
समुच्चय~$\Setabs{x}{!A(x)}$ विद्यमान है---अतः पहली दृष्टि में यह सत्य, उपयोगी
और शायद आवश्यक अभिगृहीत भी लगता है। इसे ``अनौपचारिक'' इसलिए कहते हैं क्योंकि,
जैसा अंततः पता चलता है, यह इस सिद्धांत को असंतुष्टनीय बना देती है: यदि
$!A(y)$ को~$\lnot y \in y$ लें, तो !!{sentence}
\[
\lexists[x][\lforall[y][(y \in x \liff \lnot y \in y)]]
\]
मिलता है, और यह !!{sentence} किसी भी !!{structure} में संतुष्ट नहीं होता।
\end{ex}

\begin{ex}
\emph{भागमीमांसा} में \emph{भाग-संबंध} मूलभूत संबंध है। समुच्चयों के
सिद्धांतों की तरह भाग-संबंध के भी ऐसे सिद्धांत हैं जो इस संबंध की विभिन्न, और
कभी-कभी परस्पर विरोधी, धारणाओं को अभिगृहीत करते हैं।

भागमीमांसा की भाषा में केवल एक द्विस्थानी विधेय-चिह्न~$\Obj P$ होता है, और
$\Atom{\Obj P}{x, y}$ का ``अर्थ'' है कि~$x$,~$y$ का भाग है। इस व्याख्या को
ध्यान में रखते हुए इस भाषा की किसी !!a{structure} को \emph{भाग-संरचना}
कहते हैं। निस्संदेह एक द्विस्थानी विधेय वाली हर संरचना को सचमुच यह नाम देना
उचित नहीं होगा। ``भाग-संबंध'' को निरूपित करने के लिए $\Assign{\Obj P}{M}$ को
कुछ शर्तें संतुष्ट करनी चाहिए, जिन्हें हम भाग-संबंध के सिद्धांत के अभिगृहीत
बना सकते हैं। उदाहरणतः भाग-संबंध वस्तुओं पर आंशिक क्रम है: प्रत्येक वस्तु स्वयं का एक भाग
है, भले ही \emph{अनुचित} भाग हो; कोई दो भिन्न वस्तुएँ एक-दूसरे का भाग नहीं
हो सकतीं; और किसी वस्तु के भाग का भाग स्वयं उस वस्तु का भाग है। ध्यान दें कि
इस अर्थ में ``का भाग है'', ``का उपसमुच्चय है'' से मिलता-जुलता है, पर ``का
अवयव है'' से नहीं, क्योंकि बाद वाला न स्वतुल्य है न संक्रामक।
\begin{align*}
& \lforall[x][\Atom{\Obj P}{x,x}] \\
& \lforall[x][\lforall[y][((\Part{x}{y} \land \Part{y}{x})
      \lif \eq[x][y])]] \\
& \lforall[x][\lforall[y][\lforall[z][((\Part{x}{y} \land
        \Part{y}{z}) \lif \Part{x}{z})]]]\\
\intertext{इसके अतिरिक्त, किन्हीं दो वस्तुओं का भागमीमांसात्मक योग होता है
  (ऐसी वस्तु जिसके ये दोनों वस्तुएँ भाग हैं और जो इस दृष्टि से न्यूनतम है)।}  &
\lforall[x][\lforall[y][\lexists[z][\lforall[u][(\Part{z}{u} \liff
        (\Part{x}{u} \land \Part{y}{u}))]]]]
\end{align*}
ये तत्वमीमांसकों द्वारा विचारित भाग-संबंध के मूल सिद्धांतों में से केवल कुछ हैं।
किंतु कुछ परिभाषित संबंध पहले प्रस्तुत किए बिना आगे के सिद्धांतों को
सूत्रबद्ध करना या लिखना शीघ्र ही कठिन हो जाता है। उदाहरणतः भागमीमांसा में
रुचि रखने वाले अधिकांश तत्वमीमांसक निम्न को भी वैध सिद्धांत मानते हैं: जब भी
किसी वस्तु~$x$ का उचित भाग~$y$ हो, तब उसका कोई भाग~$z$ भी होता है जिसका~$y$
के साथ कोई भाग उभयनिष्ठ नहीं होता, और~$y$ तथा~$z$ का संलयन~$x$ होता है।
\end{ex}

\end{document}
